\scanpage{56}
\mspage{12}
% archive sheet number: 55
\deleted{\unsure{\fbox{2)}|\fbox{B)}}} \inserted{(B)} \uline{Opération} de \deleted{\uline{affaiblissement} \unsure{du}} \inserted{sur}\uline{marquage}

\deleted{On peut \unsure{le|la} définir \unsure{en termes de|par les données}, \unsure{et avec} pour chaque \unsure{arête} \illegible{2} \unsure{se donner} \illegible{1} $R$ \unsure{$\widetilde{\Sigma}\setminus R$}. Soit $\deleted{\unsure{\widetilde\Sigma_*}}A^{+}_{\unsure{*|0}}\subset A_*$ l'ens}

\deleted{\unsure{des}} Soit $A^+$ l'ens des arêtes \inserted{(libres ou liées)} ``effectivement marquées'', \inserted{du $\Gamma$, i.e. telles que $r(a)>0$} \deleted{et $B$ \unsure{l'ens des}},
correspondant à des $C\in\pi_0(\Sigma)$ telles que $C\cap R\ne\unsure{\varnothing|\phi}$, soit $B^+$ l'ens des comp. connexes de $ {\unsure{\Sigma|\underline{\Sigma}}}^{+}\setminus R\cap\Sigma^{+}$ (lesquelles sont des int. ouverts), un surmarquage est défini par la donnée d'une application
% slanted
\margin{\deleted{\unsure{arêtes}} \unsure{cas des|ce des} \deleted{\unsure{repères}} \inserted{effectivement} marqués), \unsure{des} \uline{arêtes} \unsure{sur} $\Sigma$ (\uline{NB} elles se trouvent sur des cercles eff. marqués, et sont \unsure{donc} des intervalles ouverts)}
\diagram{
\begin{tikzcd}[ampersand replacement=\&]
B^{+}\arrow[r,"r'"]\&\mathbb N\ ,
\end{tikzcd}
}
et peut se décrire \deleted{par} \inserted{comme} l'opération qui \inserted{\unsure{met}} sur chaque \uline{arête} \deleted{\unsure{(ouverte) qui}} de $\Sigma$, correspondant à un $b\in B^{+}$, $r'(\deleted{\unsure{a}}b)$ nouveaux sommets de \deleted{\unsure{marquage}}\inserted{rigidification}. Cette fois-ci, \deleted{\unsure{l'opér.}} il y a des choix arbitraires --- néanmoins l'opération donne un objet de $S{\unsure{\mathscr T|\mathcal T}}D^{+}$, défini \inserted{soit $(X,\Sigma,\deleted{\unsure{R}}R_{1})$} à isom unique près --- muni bien entendu d'un ens marqué $R'\subset R_{1}$), tel que l'opération inverse de gommage de $R'$ redonne l'objet de départ $(X,\widetilde{\Sigma},R)$.
